\chapter*{भूमिका (Preface)}
\addcontentsline{toc}{chapter}{भूमिका (Preface)}
\markboth{भूमिका}{भूमिका}

इक्कीसवीं सदी का विज्ञान केवल प्रयोगशाला (laboratory) तक सीमित नहीं रह गया है।
आज हर प्रयोग, हर मापन (measurement) और हर खोज के पीछे \textbf{data (data)} और
\textbf{कृत्रिम बुद्धिमत्ता (Artificial Intelligence, AI)} की भूमिका तेज़ी से बढ़ रही है।
मौसम का पूर्वानुमान, नई दवाओं की खोज, अंतरिक्ष-यानों का संचालन, या किसी रोग की पहचान -
इन सभी में अब AI एक सहायक वैज्ञानिक (assistant scientist) की तरह कार्य करती है।

यह पुस्तक युवा पाठकों एवं विज्ञान-प्रेमियों के लिए लिखी गई है। इसका उद्देश्य
विद्यार्थियों को AI का गणित या प्रोग्रामिंग कठिन रूप में रटाना नहीं, बल्कि यह दिखाना है
कि \emph{भौतिकी (Physics), रसायन विज्ञान (Chemistry) एवं जीव विज्ञान (Biology)} में AI
का उपयोग कैसे होता है। इसलिए हर अध्याय \textbf{गतिविधि (activity)} और
\textbf{परियोजना (project)} पर आधारित है, ताकि विद्यार्थी स्वयं करके सीखें।

\medskip
\noindent\textbf{पुस्तक की भाषा-शैली:} सरल हिन्दी में व्याख्या दी गई है, परन्तु सभी
\emph{वैज्ञानिक एवं तकनीकी शब्द (scientific and technical terms)} अंग्रेज़ी में रखे गए हैं,
क्योंकि उच्च शिक्षा एवं प्रतियोगी परीक्षाओं में यही शब्दावली प्रचलित है।

\medskip
\noindent\textbf{पुस्तक में प्रयुक्त संकेत (icons/boxes):}
\begin{itemize}[leftmargin=1.4em]
  \item \textbf{गतिविधि (Activity)} - कक्षा या घर पर स्वयं करने योग्य कार्य।
  \item \textbf{परियोजना (Project)} - समूह में करने योग्य बड़ा कार्य।
  \item \textbf{क्या आप जानते हैं?} - रोचक तथ्य।
  \item \textbf{सोचिए (Think it over)} - विचार एवं चर्चा हेतु प्रश्न।
\end{itemize}

\medskip
\noindent किसी भी गतिविधि को करने के लिए महँगे उपकरण आवश्यक नहीं हैं - एक मोबाइल/कंप्यूटर,
इंटरनेट और जिज्ञासा ही पर्याप्त है। आइए, विज्ञान को AI की दृष्टि से देखना आरम्भ करें।
